\textbf{COLUNI 2019 -- Questão 36}\par\medskip Leia o texto abaixo e depois responda o que se pede:

\par

RIO — O jurista francês René Cassin não queria proteger um ou outro grupo específico de pessoas quando disse: ―Não haverá paz neste planeta enquanto os direitos humanos forem violados em alguma parte do mundo‖. Um dos autores do texto da Declaração Universal dos Direitos Humanos, adotada pela Organização das Nações Unidas (ONU) em 1948, o ganhador do Nobel da Paz de 1968 incluiu todos os Homo sapiens naquela frase célebre. Morto em 1976, aos 88 anos, Cassin seria uma ótima pessoa para invocar diante de compreensões equivocadas sobre a expressão ―direitos humanos‖, quase 70 anos depois da adoção do texto pela comunidade internacional. [...] O respeito pelos direitos de todos os humanos, explícito na declaração, é o tema do Dia Internacional da Paz, celebrado nesta sexta-feira. Num texto sobre esta data (da Proclamação da Declaração Universal dos Direitos Humanos) divulgado em seu site, a Organização das Nações Unidas para a Educação, a Ciência e a Cultura (Unesco) afirma que os direitos humanos ―são um pré-requisito para uma sociedade pacífica‖. Mas a mesma mensagem alerta para desafios no caminho dessa paz, como a desigualdade social, os conflitos gerados pelas mudanças climáticas e as visões extremistas que se espalham pelo mundo. (Disponível em: <https://oglobo.globo.com/sociedade/os-direitos-humanos-nao-sao-da-esquerda-ou-da-direita-sao-de-todos-23088573>. Acesso em: 2 out. 2018. Adaptado.)

\par

A criação da Organização das Nações Unidas aconteceu no contexto do final da Segunda Guerra Mundial e a Declaração Universal dos Direitos Humanos, proclamada em 1948, celebrou 70 anos em 2018.

\par

A partir da leitura do texto e do seus conhecimentos sobre a Declaração Universal dos Direitos Humanos, assinale a alternativa INCORRETA:\begin{enumerate}
\item Todos os seres humanos nascem livres e iguais em dignidade e direitos. São dotados de razão e
      consciência e devem agir em relação uns aos outros com espírito de fraternidade.
\item Os direitos estabelecidos nessa Declaração não encontram obstáculos e desafios no contexto do mundo
      contemporâneo.
\item Todo ser humano tem capacidade de gozar os direitos e as liberdades [...], sem distinção [...] de raça, cor,
      sexo, língua, religião, opinião política [...], origem nacional ou social.
\item Ninguém será mantido em escravidão ou servidão; [...] Ninguém será submetido à tortura, nem a
      tratamento ou castigo cruel, desumano ou degradante.\end{enumerate}
